\chapter{STL、迭代器、算法与 ranges}

\begin{introduction}
  \item 根据访问模式选择连续、节点式和关联容器
  \item 用迭代器理解算法与数据表示的解耦
  \item 使用 ranges 构造惰性、可组合的数据视图
  \item 识别复杂度、分配与迭代器失效成本
\end{introduction}

\chaptermeta{会使用 Python 列表、字典、生成器与推导式。}{ranges 视图类似惰性生成器，但受静态类型、借用范围和迭代器失效约束。}{结果异常时先检查视图底层对象寿命、迭代器失效、谓词副作用和浮点归约顺序。}

\section{先从 vector 开始}

\texttt{std::vector<T>} 的连续存储对缓存和预取友好，通常应成为序列容器默认选择。只有测量和操作需求证明频繁中间插入、稳定地址或双端增长重要时，才考虑其他结构。链表虽然单次已定位插入是常数时间，但遍历中的指针追逐和分配成本常使它输给移动一段连续内存。

\texttt{reserve} 预留容量但不改变元素数量；\texttt{resize} 改变大小并构造元素。扩容可能使所有指针、引用和迭代器失效，这是常见生命周期错误来源。

\section{算法表达意图}

\texttt{std::sort}、\texttt{std::transform}、\texttt{std::accumulate} 和查找算法比手写循环更直接表达操作类别，也让复杂度预期清楚。算法不是为了消灭所有循环；涉及多状态推进、早停或需要精细诊断的业务逻辑，命名良好的普通循环可能更可读。

迭代器把“元素位置”抽象为协议。随机访问算法要求更强迭代器能力，concepts 会在编译期拒绝不满足要求的范围。不要默认每个迭代器都支持加整数或常数时间距离。

\section{ranges 是视图，不是新容器}

视图通常不拥有元素，而是惰性描述如何遍历底层范围：

\lstinputlisting[caption={筛选买入成交并聚合名义金额}]{code/ch04/ranges_pipeline.cpp}

\texttt{filter} 不立即复制买入成交；循环拉取时才执行谓词。优点是组合和避免中间容器，代价是视图生命周期依赖底层范围，重复遍历可能重复计算。视图是否更快必须结合访问次数和编译结果测量。

\begin{pythonbridge}
生成器表达式与 ranges 视图都具有惰性，但 \cpp{} 视图的类型编码了适配器组合，并受到借用范围和迭代器失效规则约束。把临时容器接入长期保存的视图前，必须确认所有权。
\end{pythonbridge}

\section{关联查找与散列}

\texttt{std::map} 提供有序树结构和对数复杂度，\texttt{std::unordered\_map} 提供平均常数查找但有散列、桶和重哈希成本。持仓按代码查找常可从后者开始；若需要有序遍历、范围查询、稳定最坏界或确定性输出，则树结构可能更合适。

小数据集上，排序向量加二分查找经常优于节点容器，因为连续布局减少间接访问。选择数据结构应同时考虑规模、更新频率、遍历模式和延迟尾部，而不是只背复杂度表。

\section{数值归约的注意事项}

浮点加法不满足结合律，并行归约或改变遍历顺序可能改变末位结果。绩效统计应规定允许误差、稳定算法和可重复性要求。对于金额，先决定表示和舍入规则；对于收益与方差，考虑补偿求和或在线稳定算法，而不是盲目累加。

\section{常见错误}

\begin{enumerate}
  \item 在遍历 \texttt{vector} 时插入并继续使用旧迭代器。
  \item 为一次线性扫描建立昂贵的散列表。
  \item 返回引用临时容器的视图。
  \item 在谓词中加入不可见副作用，使算法结果依赖调用次数。
  \item 只比较理论复杂度，不测量真实规模下的缓存和分配成本。
\end{enumerate}

\begin{interviewcheck}
为什么 \texttt{vector} 常比 \texttt{list} 快？\texttt{reserve} 与 \texttt{resize} 有什么区别？ranges 视图拥有什么、不拥有什么？何时有序向量可能胜过哈希表？
\end{interviewcheck}

\section{练习}

\begin{enumerate}
  \item 用 ranges 筛选指定代码的买入成交，并计算加权平均价格。
  \item 对比 \texttt{vector} 线性扫描、排序后二分查找和 \texttt{unordered\_map} 在不同持仓数量下的成本。
  \item 构造一次扩容导致引用失效的最小例子，并提出两种修复。
  \item 解释并行浮点归约为何可能不逐位可重复，定义合理测试断言。
\end{enumerate}

\section{本章小结}

STL 的核心不是容器清单，而是表示、迭代协议和算法复杂度的组合。连续存储是强默认，ranges 提供非拥有惰性视图，任何性能判断都必须回到真实数据规模和访问模式。
